% Answers, scoring guides and placement decision guide for assessments.tex.
% Included by answer-key.tex (no \documentclass here).
\kcpart{A}{Assessment Answers}{Placement quiz, week-6 checkpoint and week-12 exit assessment: answers, scoring, rubrics and next steps.}

\begin{teachernote}[How to score (read once)]
\begin{itemize}[itemsep=1pt]
  \item Mark in pencil, with check marks and short comments. Never show a total to the group; talk
        about results privately and start with something the child did well.
  \item Give partial credit freely when the thinking is right but a detail is wrong (for example the
        right lines in the wrong order). The points below say where.
  \item Spelling and handwriting do not matter unless they would change what Python does
        (a missing colon or quote does matter in a \enquote{fix the bug} question).
  \item Each paper ends with a four-level scale: \textbf{Just starting}, \textbf{Getting there},
        \textbf{Got it}, \textbf{Expert}. Use the score bands to choose a level, then compare it
        with the child's own self-assessment and talk about any difference.
\end{itemize}
\end{teachernote}

%% =====================================================================
\section{Placement quiz (start of week 1)}
Total: \textbf{14 points} plus 1 bonus point. About 3 minutes per child to score.

{\small
\begin{longtable}{@{}>{\bfseries}p{1.3cm}>{\raggedright\arraybackslash}p{4.9cm}>{\raggedright\arraybackslash}p{1.2cm}>{\raggedright\arraybackslash}p{7.0cm}@{}}
\toprule
Q & Answer & Points & If it is wrong, it tells you\ldots \\ \midrule\endhead
1 & \texttt{Hip}, \texttt{Hip}, \texttt{Hip}, \texttt{Hooray!} (four lines) & 2 & 1 point for three Hips, 1 for a single Hooray. Several Hoorays means indentation is not yet clear: revisit it in week 1. \\
2 & \texttt{8} & 1 & \enquote{5} or \enquote{coins + 3}: a variable as a box that changes needs revisiting (week 1 warm-up). \\
3 & B, a square & 1 & Loops and turning angles: show the square on paper, 4 sides $\times$ 90 degrees. \\
4 & a) \texttt{Medium} \quad b) \texttt{Medium} & 2 & b) is the edge case: 10 is not bigger than 10. Wrong b) is very common and fine; it is exactly what week 8 (edge cases) is about. \\
5 & Line 3 needs a colon: \texttt{for i in range(3):} \quad Line 4: \texttt{t.forward(80)} & 2 & 1 point per bug. Missing both: the child needs practice reading code slowly; pair them with a careful navigator. \\
6 & 1, 2, 3, 4, 5 and 6 (both ends included) & 1 & Accept \enquote{1 to 6}. \enquote{1 to 5} is a common slip: \texttt{randint} includes 6, \texttt{range} does not include its end. \\
7 & a) 2 \quad b) the second one (size 100) & 2 & Functions with parameters are new-ish; week 2 builds on them. A wrong a) suggests reviewing what calling a function means. \\
8 & see below & 3 & 1 point for \texttt{for} with \texttt{range(4)}, 1 for the colon and indentation, 1 for the correct \texttt{print}. \\
Bonus & see below & +1 & Any way that prints 1--4 in order with the message. \\
\bottomrule
\end{longtable}}

\begin{answer}[Question 8 and bonus (one correct way)]
\begin{pycode}[Question 8]
for i in range(4):
    print("Go team!")
\end{pycode}
\begin{pycode}[Bonus]
for i in range(1, 5):
    print(i)
    print("Go team!")
\end{pycode}
Also accept \py|print(i + 1)| with \py|range(4)|, or \py|print(f"{i} Go team!")|.
\end{answer}

\needspace{16\baselineskip}
\subsection{Placement decision guide: choosing a path}
\begin{rubric}[Placement: score bands]
\renewcommand{\arraystretch}{1.3}
\begin{tabularx}{\linewidth}{@{}>{\bfseries}p{2.2cm}>{\bfseries}p{2.6cm}X@{}}
\toprule
Score & Path & What it means in weeks 1--5 \\ \midrule
0--6 & Support path & Use the support path in every lesson plan: cheat sheet beside the keyboard, fill-in-the-blank code, the child drives first while a partner navigates. Spend 5 minutes of week 1 and 2 warm-ups on \texttt{for} loops and \texttt{if}. If most of the group is here, consider using B2 (Generative Art Jam) early as a Level 1 refresher. \\
7--11 & Standard path & Follow the lesson plans as written. Offer the extension task when a child finishes early. \\
12--14 (+bonus) & Extension path & Offer the extension path and stretch challenge each week. Give this child the code reviewer role often and ask them to explain, not to type for others. \\
\bottomrule
\end{tabularx}
\end{rubric}

\begin{tip}[Look at the pattern, not only the total]
\begin{itemize}[itemsep=1pt]
  \item \textbf{Q1 and Q8 wrong} (loops): put this child on the support path for weeks 1--4
        even if the total is in the standard band. Loops underpin \texttt{while} (week 4).
  \item \textbf{Q7 wrong} (functions): give extra time to the function machine in week 2.
  \item \textbf{Q5 wrong but everything else right}: usually rushing, not lack of understanding.
        Teach \enquote{read it twice, slowly} and the rubber-duck habit.
  \item \textbf{Low score but strong when you ask the questions out loud}: reading, not coding, is
        the difficulty. Read questions aloud from now on and offer the dyslexia-friendly edition.
  \item A child new to the group with no Level 1 experience: support path, and ask them to try a few
        Level 1 week 7--10 exercises at home in the first month.
  \item Paths are not labels. Write the path in pencil on the progress tracker and look again after
        the week-6 checkpoint.
\end{itemize}
\end{tip}

\clearpage
\section{Mid-level checkpoint (week 6)}
Total: \textbf{18 points}.

{\small
\begin{longtable}{@{}>{\bfseries}p{0.6cm}>{\raggedright\arraybackslash}p{6.4cm}>{\raggedright\arraybackslash}p{1.2cm}>{\raggedright\arraybackslash}p{6.2cm}@{}}
\toprule
Q & Answer & Points & Notes \\ \midrule\endhead
1 & \texttt{42}: int \quad \texttt{3.5}: float \quad \texttt{"42"}: str \quad \texttt{True}: bool & 2 & 4 right = 2 points, 2--3 right = 1. Mixing up \texttt{42} and \texttt{"42"} is the key misconception (week 1 bug hunt: \texttt{"5" + 5}). \\
2 & \texttt{Kai has 6 stars} & 1 & The math inside \texttt{\{ \}} is worked out first. \\
3 & \texttt{9} then \texttt{4} & 2 & 1 point each. \texttt{double(double(1))}: the inner call runs first and gives 2. \\
4 & \texttt{4} then \texttt{sky blue} & 2 & \enquote{orange} for the second line means the child still counts from 1; revisit the human list. \\
5 & Trace: 10, 7, 4, 1, $-2$. It prints \texttt{-2} & 2 & 1 point for the trace, 1 for $-2$. Many children stop at 1: the loop checks the condition only at the top. \\
6 & \texttt{30} & 1 & \\
7 & \texttt{area} prints instead of returning, so it gives back \texttt{None}. Fix: \texttt{return width * height}. The program then prints \texttt{21}. & 2 & 1 point for spotting print vs return, 1 for the fix. Before crashing, the buggy program prints 20, which confuses many children. \\
8 & \texttt{print(pets["cat"])} (keys are case-sensitive). Prints \texttt{2}. & 1 & \\
9 & see below & 3 & 1 point for looping over the list with an \texttt{if}, 1 for counting, 1 for \texttt{return} (not \texttt{print}). \\
10 & At the \textbf{last line}: it names the error and says what went wrong. The line just above shows the line number in our code. & 2 & 1 point for \enquote{the bottom}, 1 for a reason. \\
\bottomrule
\end{longtable}}

\begin{answer}[Question 9 (one correct way)]
\begin{pycode}
def count_stars(items):
    count = 0
    for item in items:
        if item == "star":
            count = count + 1
    return count
\end{pycode}
Also correct: \py|return items.count("star")|. Check that the function \emph{returns}: if it
prints, give 2 of 3 points and ask \enquote{How could another function use your answer?}
\end{answer}

\begin{rubric}[Checkpoint: score bands and next steps]
\renewcommand{\arraystretch}{1.3}
\begin{tabularx}{\linewidth}{@{}>{\bfseries}p{1.6cm}>{\bfseries}p{2.4cm}X@{}}
\toprule
Score & Level & Next steps for weeks 7--12 \\ \midrule
0--6 & Just starting & Support path. Use B1 or B2 soon as a calm review session. In the project, start from the starter template and aim for milestone M2. \\
7--11 & Getting there & Standard path, with the support task for the goal they found hardest (see the goal map below). \\
12--15 & Got it & Standard path; offer extensions. \\
16--18 & Expert & Extension path. In the project, suggest a separate logic file like the Star Catcher reference. \\
\bottomrule
\end{tabularx}
\end{rubric}

\textbf{Goal map} (which questions show which goal): L2.G1 types and f-strings: Q1, Q2.
L2.G2 while loops: Q5. L2.G3 functions and return: Q3, Q7, Q9. L2.G4 lists and dictionaries:
Q4, Q6, Q8, Q9. L2.G6 debugging: Q7, Q8, Q10. Record a level for each goal on the group page
of the progress tracker.

\clearpage
\section{Exit assessment (week 12)}
Total: \textbf{20 points}. Question 9 is about the child's own project: score it after the showcase
if that is easier.

{\small
\begin{longtable}{@{}>{\bfseries}p{0.6cm}>{\raggedright\arraybackslash}p{6.6cm}>{\raggedright\arraybackslash}p{1.2cm}>{\raggedright\arraybackslash}p{6.0cm}@{}}
\toprule
Q & Answer & Points & Notes \\ \midrule\endhead
1 & \texttt{2 cat} & 2 & 1 point for 2 (two lines), 1 for \texttt{cat}. \\
2 & Only \texttt{fish}: opening with \texttt{"w"} empties the file before writing. & 1 & \\
3 & The file does not exist yet the first time. See the fix below. & 2 & 1 point for the reason, 1 for \texttt{try}/\texttt{except FileNotFoundError} giving an empty list. Also accept checking that the file exists first. \\
4 & a) pass. b) The \textbf{test} is wrong: clamp keeps the value at least \texttt{low}, so the right answer is 0. c) e.g. \texttt{assert clamp(-3, 0, 10) == 0} or \texttt{assert clamp(10, 0, 10) == 10} & 3 & 1 point each. An edge case is at or past a limit (the walls, 0, an empty list). \\
5 & \texttt{screen.onkey(jump, "space")}: no brackets, so turtle can call \texttt{jump} later, each time space is pressed. \texttt{listen()} makes the window pay attention to the keyboard. & 2 & 1 point each. \\
6 & a) 5 times \quad b) 5 & 2 & 1 point each. The first call is the one \enquote{somewhere else}; then 4 more by the timer. \\
7 & It checks whether the basket and the star are closer than 30 steps (pixels), i.e. touching: the star is caught (a collision). & 1 & \\
8 & see below & 4 & 1 point each: handles the empty list, finds the biggest, uses \texttt{return}, one sensible test. \\
9 & Own project & 3 & 1 point each for parameters (what goes in), return value (what comes out), and how it was tested. \\
\bottomrule
\end{longtable}}

\begin{answer}[Question 3]
\begin{pycode}
try:
    with open("scores.txt") as file:
        scores = file.read().split()
except FileNotFoundError:
    scores = []
\end{pycode}
\end{answer}

\begin{answer}[Question 8 (one correct way)]
\begin{pycode}
def best_score(scores):
    if len(scores) == 0:
        return 0
    best = scores[0]
    for score in scores:
        if score > best:
            best = score
    return best

def test_best_score():
    assert best_score([3, 9, 4]) == 9
    assert best_score([]) == 0
\end{pycode}
Also correct: \py|return max(scores)| after the empty check. A test without the empty list is
fine; praise it if the child tests the edge case too.
\end{answer}

\begin{rubric}[Exit: score bands]
\renewcommand{\arraystretch}{1.3}
\begin{tabularx}{\linewidth}{@{}>{\bfseries}p{1.6cm}>{\bfseries}p{2.4cm}X@{}}
\toprule
Score & Level & What it means \\ \midrule
0--7 & Just starting & Has met many Level 2 ideas but needs more practice before Level 3. Repeat the project with a new idea, or redo weeks 7--9 at a gentle pace. \\
8--12 & Getting there & Ready for Level 3 with support: review files and tests in the first Level 3 weeks. \\
13--16 & Got it & Ready for Level 3. \\
17--20 & Expert & Ready for Level 3; offer the extension path from the start. \\
\bottomrule
\end{tabularx}
\end{rubric}

\textbf{Goal map:} L2.G5 files: Q1, Q2, Q3. L2.G6 tests and debugging: Q4, Q5, Q8. L2.G7
events and games: Q5, Q6, Q7. L2.G3 and L2.G4 functions and lists: Q8. L2.G8 explaining: Q9.
Combine with the project rubric (Project Teacher Guide) to fill in the week-12 column of the
progress tracker.

\begin{tip}[Writing the certificate]
On each certificate, add one specific sentence about something the child did well: \enquote{Built
a dragon-feeding game with a saved high-score table} or \enquote{Wrote six tests and found a bug with
one of them}. Specific praise means far more than \enquote{well done}.
\end{tip}
